% =====================================================================
%  บทที่ 5  สรุปผล อภิปรายผล และข้อเสนอแนะ
% =====================================================================
\chapter{สรุปผล อภิปรายผล และข้อเสนอแนะ}
\label{ch:conclusion}

\section{สรุปผลการวิจัย}

การวิจัยนี้ได้พัฒนาระบบผู้ช่วยสอนด้วยปัญญาประดิษฐ์เชิงสร้างสรรค์สำหรับรายวิชาการเขียนโปรแกรมเบื้องต้น ผลการทดลองพบว่าผู้เรียนที่ใช้ระบบมีผลสัมฤทธิ์ทางการเรียนสูงกว่าผู้เรียนที่ไม่ได้ใช้ระบบอย่างมีนัยสำคัญทางสถิติ (บทที่~\ref{ch:results}) และมีความพึงพอใจต่อระบบในระดับมากที่สุด

\section{อภิปรายผล}

ผลการวิจัยสอดคล้องกับแนวคิดที่ว่า การให้ข้อมูลป้อนกลับทันทีช่วยให้ผู้เรียนเข้าใจข้อผิดพลาดของตนเองได้ดีขึ้น นอกจากนี้ การใช้เทคนิค RAG ทำให้คำตอบของระบบอ้างอิงเนื้อหาของรายวิชา ลดปัญหาการสร้างข้อมูลที่ไม่ถูกต้อง (hallucination) ของแบบจำลองภาษา

\section{ข้อเสนอแนะ}

\subsection{ข้อเสนอแนะในการนำผลการวิจัยไปใช้}
\begin{enumerate}
  \item ผู้สอนควรกำหนดแนวทางการใช้งานระบบให้ชัดเจน เพื่อป้องกันการคัดลอกคำตอบโดยไม่ทำความเข้าใจ
  \item ควรปรับปรุงเอกสารประกอบการสอนในฐานข้อมูลให้เป็นปัจจุบันอยู่เสมอ
\end{enumerate}

\subsection{ข้อเสนอแนะในการวิจัยครั้งต่อไป}
\begin{enumerate}
  \item ควรศึกษาผลของระบบในระยะยาวตลอดภาคการศึกษา
  \item ควรขยายการใช้งานไปยังรายวิชาอื่น เช่น โครงสร้างข้อมูลและอัลกอริทึม
\end{enumerate}
